% =====================================================================
%  CAPÍTULO 8 — CONCLUSÕES E TRABALHO FUTURO
% =====================================================================
\chapter{Conclusões e Trabalho Futuro}
\label{cap:conclusoes}

\section{Síntese e Cumprimento dos Objetivos}
\label{sec:sintese}

Este projeto demonstrou a viabilidade de um sistema completo que
transforma um arquivo familiar heterogéneo e desorganizado em narrativas
escritas e vídeos documentais coerentes, através de um \emph{pipeline}
modular de quatro estágios assente em inteligência artificial generativa.

Revendo os objetivos definidos no \Cref{sec:objetivos}: a ingestão
multimodal (O1), a organização temporal (O2), a geração narrativa (O3) e
a produção de vídeo (O4) foram \textbf{implementadas e validadas}; a
disponibilização através de uma aplicação Web moderna e implantada na
nuvem (O5) foi \textbf{concretizada} (FastAPI + React + Supabase + Render +
Vercel); e a premissa de privacidade e baixo custo (O6) foi
\textbf{assegurada} pela estratégia de IA híbrida e pelo uso de
ferramentas de código aberto. Em suma, os objetivos centrais do projeto
foram alcançados.

Para além dos objetivos definidos, o trabalho concretizou vários avanços
que elevam o resultado: a \textbf{narrativa segmentada em cenas}, que
sincroniza cada fotografia com a sua narração no vídeo (transformando o
\emph{slideshow} num verdadeiro documentário); um \textbf{executor de
tarefas \emph{in-process}} que substitui um \emph{worker} Celery pago no
plano gratuito (usado, por exemplo, na análise diferida das fotografias),
com limpeza de tarefas órfãs no arranque; uma \textbf{árvore familiar
interativa} (React~Flow) com pesquisa, destaque de linhagem, exportação em
imagem e edição/adição de familiares; legendas de vídeo como \textbf{faixa
WebVTT} sincronizada palavra-a-palavra, com escolha de voz; e ferramentas
de \textbf{curadoria de histórias} (favoritos, pesquisa/filtro,
regeneração orientada por \emph{feedback} e exportação para PDF). A
interface está totalmente \textbf{bilingue} (português europeu e inglês),
reforçada por um pós-processador determinístico que assegura a variante
europeia mesmo quando o modelo recai no português do Brasil.

Mais do que a integração de modelos, o trabalho destacou-se pela
\textbf{engenharia de robustez}: a degradação graciosa do RAG, os
\emph{fallbacks} de IA, o \emph{sweep} de tarefas órfãs, a execução
assíncrona \emph{in-process} e a recuperação de \emph{cold starts}
permitem ao sistema operar de forma fiável quer numa máquina local com
todas as dependências, quer num ambiente de nuvem restrito. A migração de
um protótipo \emph{local-first} para uma arquitetura multi-inquilino,
preservando a opção de inferência local, ilustra um compromisso maduro
entre acessibilidade, privacidade e escalabilidade.

\section{Limitações}
\label{sec:limitacoes}

Em honestidade técnica, reconhecem-se as seguintes limitações da versão
atual:

\begin{itemize}[leftmargin=*]
  \item \textbf{Avaliação quantitativa incompleta.} Embora os dados de
  base sejam persistidos, falta ainda um painel consolidado de métricas
  (tempos de geração, taxa de sucesso por \emph{template}, qualidade do
  RAG) e um protocolo de avaliação humana sistemático.
  \item \textbf{RAG inativo em produção.} O passo de \emph{retrieval} já
  está ligado à geração (\Cref{sub:geracao_narrativa}), mas no \emph{deploy}
  gratuito a ausência de ChromaDB força o modo \emph{stub} — só em ambiente
  local o \emph{retrieval} fica plenamente ativo. A migração para
  \emph{pgvector} (\Cref{sec:trabalho_futuro}) resolveria isto.
  \item \textbf{Sincronização de vídeo aproximada.} A duração de cada
  fotografia segue a narração da sua cena, mas pequenas derivas podem
  acumular-se em histórias longas; alinhar a troca de fotografias ao nível
  da palavra --- aproveitando os carimbos temporais por palavra já
  capturados para as legendas --- seria mais preciso.
  \item \textbf{Sem identificação de pessoas.} O reconhecimento facial
  (que tornaria as narrativas menos genéricas) não foi implementado.
  \item \textbf{Cobertura de testes do \emph{frontend}.} A automação de
  testes incide sobre o \emph{backend}; faltam testes ponta-a-ponta
  (E2E) da interface.
  \item \textbf{Acessibilidade não auditada formalmente.} Existem boas
  práticas, mas falta uma auditoria \textsc{wcag}~AA.
  \item \textbf{Geração de vídeo limitada pela memória no \emph{deploy}
  gratuito.} A montagem com MoviePy/FFmpeg a 720p excede os 512~MB de RAM do
  plano gratuito do Render, que termina a instância (\emph{out-of-memory}) a
  meio do \emph{render}. A solução adotada é uma \textbf{política de geração
  exclusivamente local}: na nuvem, o \emph{backend} \textbf{recusa} o
  \emph{render} (\texttt{VIDEO\_LOCAL\_ONLY}), evitando o
  \emph{out-of-memory}; os vídeos são pré-gerados em ambiente local (a 720p,
  sem limite de memória), escritos no \emph{mesmo} \emph{Storage} (Supabase)
  e reproduzidos no \emph{site} sem qualquer \emph{render} em tempo real ---
  a reprodução é leve (apenas servir o ficheiro) e nunca é afetada. A
  resolução e os \emph{fps} permanecem parametrizáveis por ambiente como
  ajuste adicional.
  \item \textbf{Quota de visão limitada no nível gratuito.} A análise visual
  das fotografias (M1) depende do Gemini, cujo nível gratuito permite apenas
  poucas dezenas de pedidos por dia; a geração de texto (M3) não é afetada por
  recorrer ao Groq como alternativa. A re-análise manual de fotografias é, por
  isso, parcimoniosa (só processa fotografias ainda sem descrição e confirma o
  custo em pedidos antes de avançar).
\end{itemize}

\section{Trabalho Futuro}
\label{sec:trabalho_futuro}

Identificam-se as seguintes direções de trabalho futuro, por ordem de
impacto esperado:

\begin{enumerate}[leftmargin=*]
  \item \textbf{Implantação local privada.} Executar o sistema completo numa
  máquina dedicada (com \textsc{gpu}) a correr um \textsc{llm} local via
  Ollama, eliminando a dependência de modelos na nuvem. \emph{Exemplo:} uma
  família gera as histórias no seu próprio servidor, sem que os dados saiam de
  casa.

  \item \textbf{Reconhecimento facial.} Identificar automaticamente as
  pessoas nas fotografias, com confirmação do utilizador, para tornar as
  narrativas mais específicas. \emph{Exemplo:} construir a linha temporal de
  uma pessoa a partir das fotografias em que é reconhecida.

  \item \textbf{Consulta conversacional do arquivo.} Permitir perguntas em
  linguagem natural sobre o arquivo, respondidas pelo subsistema \textsc{rag}.
  \emph{Exemplo:} responder a ``que pessoas aparecem nas fotografias de
  1985?''.

  \item \textbf{Mapa geo-temporal.} Posicionar as fotografias num mapa a
  partir das coordenadas \textsc{gps} do \textsc{exif}. \emph{Exemplo:}
  visualizar o percurso geográfico da família ao longo dos anos.

  \item \textbf{Ingestão de áudio e vídeo.} Estender o módulo de ingestão a
  vídeos e gravações de voz, transcritos automaticamente (e.g.\ com o
  \emph{Whisper}). \emph{Exemplo:} incorporar o relato gravado de um familiar
  na narrativa.

  \item \textbf{Clonagem de voz.} Narrar com a voz de um familiar a partir de
  uma amostra de áudio, mediante consentimento. \emph{Exemplo:} um
  documentário narrado na voz do próprio.

  \item \textbf{Avaliação formal.} Validar a qualidade das narrativas com um
  painel humano e métricas automáticas. \emph{Exemplo:} pontuar a coerência e
  a fidelidade factual numa amostra de histórias geradas.
\end{enumerate}

\section{Considerações Finais}
\label{sec:final}

O resultado deste projeto é uma plataforma funcional, segura e
estendível, que materializa uma ideia com valor humano evidente:
preservar e dar voz à memória familiar. Demonstra, em concreto, que é
hoje possível construir — com ferramentas de código aberto e a custo
essencialmente nulo — um sistema de IA generativa multimodal que
transforma um arquivo passivo num documentário narrado. O caminho de
evolução traçado, rumo a uma avaliação quantitativa rigorosa e a
funcionalidades de colaboração, confirma o potencial do trabalho para além
do âmbito desta licenciatura.
